\documentclass[12pt]{article}
\usepackage[margin=1in]{geometry}
\usepackage{cite}
\usepackage{url}

\title{Annotated Bibliography:\\Neural Methods for Synthesizer and Effect Parameter Matching}
\author{Music 423 Research Literature}
\date{\today}

\begin{document}

\maketitle

\section*{Overview}
This annotated bibliography covers neural and computational methods for synthesizer parameter estimation and sound matching from 2002 to 2025, including approaches based on recurrent neural networks, differentiable DSP, normalizing flows, quality-diversity optimization, and audio spectrogram transformers.

\section{Early Work (2002--2017)}

\subsection*{Alvin \& Chowning (2002) \cite{Alvin2002}}
Develops a recurrent neural network model for estimating nonlinear parameters of bowed string synthesis, demonstrating early application of neural networks to physical modeling parameter estimation.

\subsection*{van den Doel \& Pai (2007) \cite{vandenDoel2007}}
Proposes wav2shape, a method to infer rigid-body dynamics parameters from recorded sounds, enabling sound-driven physical simulation and synthesis.

\subsection*{Hsu \& Glass (2010) \cite{Hsu2010}}
Applies generative models to synthesize bowed string sounds, focusing on statistical learning approaches for expressive instrument synthesis.

\subsection*{Itoyama \& Okuno (2014) \cite{ItoyamaOkuno2014}}
Introduces a Bayesian nonparametric framework for estimating sound synthesis model parameters, providing principled uncertainty quantification in parameter estimation.

\subsection*{Barkan et al.\ (2017) \cite{Barkan2017}}
Presents DeepSynthPE, using learned reverse controls to invert synthesizers by training neural networks to predict parameter configurations from audio features.

\subsection*{Yee-King, Fedden \& d'Inverno (2018) \cite{YeeKing2018}}
Explores automatic programming of VST synthesizers using deep networks combined with evolutionary algorithms, demonstrating practical applications in sound design automation.

\section{Differentiable DSP Era (2019--2022)}

\subsection*{Barkan et al.\ (2019) \cite{Barkan2019}}
Introduces InverSynth for deep estimation of discrete synthesizer parameter configurations from audio signals using convolutional neural networks trained on quantized parameter spaces.

\subsection*{Esling et al.\ (2019) \cite{Esling2019}}
Proposes FlowSynth, using normalizing flows to learn invertible mappings between audio perception space and synthesizer parameter space, enabling universal audio synthesizer control with continuous semantic manipulation.

\subsection*{Engel et al.\ (2020) \cite{engel2020ddsp}}
Introduces DDSP (Differentiable Digital Signal Processing), combining traditional signal processing with deep learning through differentiable audio components, enabling end-to-end training and high-quality audio synthesis.

\subsection*{Mart\'{\i}nez Ram\'{\i}rez et al.\ (2021) \cite{Martinez2021}}
Presents DeepAFx for differentiable signal processing with black-box audio effects, learning to emulate and control audio effects without access to internal parameters.

\subsection*{Masuda \& Saito (2021) \cite{MasudaSaito2021}}
Applies differentiable DSP to synthesizer sound matching, demonstrating improved parameter estimation and audio reconstruction compared to non-differentiable approaches.

\subsection*{Shier (2021) \cite{Shier2021}}
PhD thesis exploring synthesizer sound matching using differentiable DSP techniques, providing comprehensive treatment of gradient-based optimization for sound synthesis parameter estimation.

\subsection*{Turian et al.\ (2021) \cite{Turian2021}}
Introduces TorchSynth for generating one billion audio sounds at a time, enabling massive-scale dataset creation for training sound matching models.

\subsection*{Andreux \& Mallat (2022) \cite{Andreux2022}}
Develops differentiable time-frequency scattering for audio synthesis, combining wavelet scattering transforms with neural networks for improved audio representation learning.

\subsection*{Chen et al.\ (2022) \cite{Chen2022}}
Presents Sound2Synth for interpreting sounds via FM synthesizer parameter estimation, focusing on frequency modulation synthesis with specialized neural architectures.

\subsection*{Steinmetz, Bryan \& Reiss (2022) \cite{Steinmetz2022}}
Introduces DeepAFx-ST for style transfer of audio effects using differentiable signal processing, enabling transfer of processing characteristics between audio examples.

\section{Quality-Diversity and Advanced Methods (2023--2024)}

\subsection*{Hafner et al.\ (2023) \cite{Hafner2023}}
Proposes mesostructures for sound synthesis, introducing intermediate representations between low-level parameters and high-level perceptual features for intuitive sound design.

\subsection*{Han, Lostanlen \& Lagrange (2023) \cite{han2023perceptualneuralphysicalsoundmatching}}
Develops perceptual-neural-physical sound matching combining plug-and-play priors with DDSP, treating synthesis parameter estimation as an inverse problem with learned regularization.

\subsection*{Han et al.\ (2023) \cite{HanInverse2023}}
Formulates synthesizer sound matching as an inverse problem using plug-and-play methods, enabling flexible incorporation of priors without retraining synthesis models.

\subsection*{Masuda \& Saito (2023) \cite{MasudaSaito2023}}
Introduces quality-diversity optimization with novelty search for synthesizer sound matching, generating diverse high-quality parameter sets that match target sounds.

\subsection*{Masuda \& Saito (2023) \cite{MasudaSaito2023improved}}
Improves semi-supervised differentiable synthesizer sound matching for practical applications, addressing generalization and training efficiency challenges.

\subsection*{Diaz \& Hayes (2023) \cite{DiazHayes2023}}
Applies differentiable DDSP to rigid-body modal synthesis, enabling gradient-based optimization of physical modeling parameters for impact sounds.

\subsection*{Nishikimi et al.\ (2023) \cite{Nishikimi2023}}
Proposes white-box parameter search for synthesizer sound matching, leveraging known synthesizer architectures to improve optimization efficiency.

\subsection*{Yang \& Chen (2023) \cite{Yang2023}}
Develops white-box FM synthesizer parameter estimation methods, exploiting the mathematical structure of frequency modulation for more accurate parameter recovery.

\subsection*{Uzrad et al.\ (2024) \cite{Uzrad2024}}
Introduces DiffMoog, a differentiable modular synthesizer inspired by Robert Moog's 1964 design, integrating subtractive, FM, and additive synthesis. Features a 2D matrix architecture where cells (rows=channels, cols=layers) host modules (Oscillator, LFO, FM Oscillator, Filter, ADSR, Mix, Tremolo) with configurable signal routing. Proposes a novel \textbf{signal-chain loss} that evaluates spectral distance at all stages of the synthesis chain, not just the final output. Key finding: Wasserstein distance along the \emph{time axis} (not frequency) significantly improves frequency estimation. Honestly reports that frequency estimation via spectral loss remains intrinsically challenging and FM chains often fail to converge due to highly non-convex loss surfaces. Open-source at \texttt{github.com/aisynth/diffmoog}.

\subsection*{Bruford, Bland \& Nercessian (2024) \cite{Bruford2024}}
Applies Audio Spectrogram Transformers (AST) to synthesizer sound matching, demonstrating superior performance over CNNs and MLPs through attention mechanisms.

\subsection*{Han, Lostanlen \& Lagrange (2024) \cite{10508449}}
Learns to solve inverse problems for perceptual sound matching, developing end-to-end trained systems that combine perception models with synthesis optimization.

\section{Recent Advances (2025)}

\subsection*{Han, Lostanlen \& Lagrange (2025) \cite{Han2025eusipco}}
Investigates gradient clipping strategies for improved training stability in DDSP models, addressing optimization challenges in differentiable synthesis.

\subsection*{Han (2025) \cite{HanT2025}}
Reviews recent plug-and-play methods for audio inverse problems, surveying advances in prior-based optimization for sound synthesis parameter estimation.

\subsection*{Pfalz \& Grill (2025) \cite{Pfalz2025metrics}}
Analyzes sound similarity metrics for synthesizer sound matching, comparing perceptual distance measures and their effectiveness for optimization objectives.

\subsection*{Hayes (2025) \cite{Hayes2025}}
Provides comprehensive survey of differentiable audio processing and synthesis, reviewing architectures, training strategies, and applications across audio domains.

\subsection*{Engel et al.\ (2025) \cite{Engel2025modulation}}
Proposes modulation discovery for DDSP, automatically learning modulation structures and parameter interactions in differentiable synthesis models.

\bibliographystyle{plain}
\bibliography{References}

\end{document}
